% !TEX TS-program = xelatex
% !TEX encoding = UTF-8 Unicode
\documentclass[11pt]{article}
\usepackage[margin=1in]{geometry}

\usepackage{xcolor}
\usepackage{hyperref}
\hypersetup{
  colorlinks=true,
  citecolor=violet
}

\usepackage[style=apa]{biblatex}
\addbibresource{references.bib}

\usepackage{fontspec} 
\setmainfont{Charis SIL}


\title{Thesis Pre-proposal: Sameer}
\author{Wren Arbor}
\date{August 15th, 2025}

\begin{document}
\maketitle
If I were to be assigned Sameer as my thesis advisor, I would investigate the research question: Is difficulty aquiring secondary tone pronunciation characteristics, such as voice quality, a primary cause of Mandarin second-language learners' difficulty with T2/T3 contrast?  Does this differ depending on the phonology of their first language(s)?  Although this would not be directly tied to any work I've done in Sameer's classes (or other linguistics classes), I would certainly rely on some of the skills and material taught in Phonetics and MDA for this thesis, especially for methodology. I landed on this question after a brief initial literature search into the area of Mandarin tone aquisition. 

In this search, I found that this area consists of two main threads: Mandarin first-language tone aquisition among children (eg. \cite{rhee_going_2021}), which investigates how children's tonal systems develop over their language development, and Mandarin second-language tone aquisition among young adult students. In the first-language literature, \textcite{rhee_going_2021} point out that although phonetics research broadly recognizes the multimodal nature of Mandarin tones beyond f0 contour -- primarily voice quality -- that this aspect of the Mandarin tone system has gone mostly unconsidered in work on tone aquisition. They find that these non-f0 contrasts are aquired late in the first-language process, achieving adult-level production 2-4 years after the f0 element. 

In the second-language literature, two issues seem to be prominent: the question of the influence of the first language (primarily English L1 vs. Cantonese L1), and the question of the almost universal difficulty in distinguishing T2 and T3 (eg. \cite{hao_second_2012, meckelborg_acquisition_2024}). Although \textcite{hao_second_2012} and \textcite{meckelborg_acquisition_2024} both find that learners have the most difficulty with T2/T3 regardless of first language(s), Hao finds that students with Cantonese first language do not perform better at tone identification, mimicry, or production than those with English first language, whereas Meckelborg et al. find (the perhaps more expected result) that the Cantonese-experienced students perform more accurately in production.

Investigation of second-language learning of non-f0 tone characteristics may be illuminating in resolving this connflict. It is possible that mastery of non-f0 tone characteristics (which may come later in the learning process as it does for L1 aqcuirers in \cite{rhee_going_2021}) is of primary importance to the mastery of T2/T3 contrasts, and that English and Cantonese speakers are on similar ground here due to lack of non-f0 tone characteristics in the Cantonese tone system (as reported by \cite{mok_assessing_2019}; cited in \cite{rose_variation_2020}). This would be evidence that Cantonese language background does not have material benefit over English language background in the aquisition of T2/T3 contrast.  It is also possible that unlike for first-language aquiring children, second-language students pick up on non-f0 tone characteristics at a similar or faster rate than f0 characteristics, and that f0 similarity of T2/T3 is the main cause of students' difficulties; this would suggest that Cantonese-background students may be advantaged by their greater f0 awareness compared to English-background students.

Methodologically, I would be comparing students' production and perception of non-f0 tone characteristics across the variables of time studying Chinese and English-background vs. Chinese-background, mainly focused on T2/T3 and creak. In production,  I would likely follow the design in \textcite{hao_second_2012} with mimicry and reading tasks, and perform auditory and acoustic analysis to evaluate the mastery of non-f0 tone characteristics. In perception, it would likely be necessary to isolate the non-f0 elements of the T2/T3 contrast using pitch correction software, and evaluate students' accuracy at tone identification on these pitch-constant stimuli.

This proposal is my number 2 pick; I find it interesting, but a little bit less exciting than my sociolinguistic idea.

\printbibliography

\end{document}